\documentclass[11pt]{article}
\usepackage[T1]{fontenc}
\usepackage[margin=1in]{geometry}
\usepackage{amsmath,amssymb}
\usepackage{enumitem}
\usepackage{hyperref}
\hypersetup{colorlinks=true,linkcolor=black,urlcolor=blue}

\begin{document}

\begin{center}
{\LARGE Coherence Formula Sheet}\\[0.5em]
{\large PolarH10 RR-derived coherence}
\end{center}

\section*{Data path}

\begin{itemize}[leftmargin=1.5em]
\item Input stream: accepted RR intervals from the Heart Rate Measurement characteristic.
\item Physical H10: those RR intervals are device-derived from ECG beat timing inside the strap.
\item Synthetic transport: the coherence solve still uses RR, while synthetic PMD ECG is emitted separately so the waveform view stays aligned to the same beat schedule.
\end{itemize}

\section*{Fixed spectral model}

The app keeps the paper-defined constants fixed:

\begin{itemize}[leftmargin=1.5em]
\item peak search band: $0.04$--$0.26$ Hz
\item peak integration window: $0.030$ Hz centered on the dominant peak
\item total-power band: $0.0033$--$0.4$ Hz
\end{itemize}

These formulas operate on the RR-derived tachogram after the irregular
beat-interval stream has been resampled into a continuous time series. In the
notation below, $f$ is frequency in hertz, $PSD(f)$ is the power spectral
density at that frequency, $f_{peak}$ is the strongest oscillation in the
allowed coherence band, $P_{peak}$ is the power concentrated around that
oscillation, and $P_{total}$ is the broader band power used as the paper's
reference total.

\[
f_{peak} = \arg\max_{f \in [0.04, 0.26]} PSD(f)
\]

This step finds the dominant rhythm inside the coherence search band instead of
across the entire spectrum.

\[
P_{peak} = \int_{f_{peak} - 0.015}^{f_{peak} + 0.015} PSD(f)\,df
\]

This integrates a narrow window around the dominant peak so the metric rewards
concentrated narrowband power rather than diffuse nearby energy.

\[
P_{total} = \int_{0.0033}^{0.4} PSD(f)\,df
\]

This is the total reference power across the wider HRV band used by the cited
method.

\[
\text{paper coherence ratio} =
\left(\frac{P_{peak}}{P_{total} - P_{peak}}\right)^2
\]

This ratio rises when one narrow resonance-like peak dominates the spectrum and
falls when power is spread across the rest of the band.

\section*{UI-facing normalized score}

The app also publishes a bounded score for the headline readout and summary charts:

\[
\text{normalized coherence}_{0..1} =
\operatorname{clamp}\left(\frac{P_{peak}}{P_{total}}, 0, 1\right)
\]

Here the app is simply reporting the fraction of tracked power that sits in the
dominant peak window, then clamping that fraction to a display-friendly range.

This normalized score is an application-level presentation choice. It is not the
McCraty paper ratio itself.

\section*{Runtime choices around the paper formula}

The cited source fixes the spectral constants above, but it does not define the
operator-side runtime details used by this app. The current implementation uses:

\begin{itemize}[leftmargin=1.5em]
\item a rolling RR window
\item cubic-spline resampling of the RR tachogram
\item a 128-point Hann-windowed spectrum solve
\item a smoothed headline $0..1$ display value
\item a separate confidence score based on RR count and window coverage
\end{itemize}

Those runtime choices describe how the app gets from incoming RR intervals to
$PSD(f)$ and to a stable operator readout. They affect warmup, smoothness, and
confidence behavior, but not the core paper ratio formula shown above.

\section*{Alignment status}

The implementation is aligned to the cited reference at the formula level for
the peak search band, peak-window width, total-band limits, and paper coherence
ratio. The display normalization, smoothing, and confidence telemetry are
explicit application-level additions.

\section*{References}

\begin{itemize}[leftmargin=1.5em]
\item McCraty, Atkinson, Tomasino, and Bradley, \emph{The Coherent Heart} (2006).
\item \url{https://mesmerprism.github.io/PolarH10/reference/coherence-workflow.html}
\item \url{https://mesmerprism.github.io/PolarH10/reference/references.html}
\end{itemize}

\end{document}
